% ===== PRÉAMBULE CANONIQUE — à coller VERBATIM en tête de CHAQUE .tex =====
\documentclass[11pt,a4paper]{article}
\usepackage[utf8]{inputenc}\usepackage[T1]{fontenc}\usepackage{lmodern}
\usepackage[french]{babel}
\usepackage{amsmath,amssymb,mathtools}\usepackage{icomma}
\usepackage{siunitx}\sisetup{locale=FR}  % \qty{}{}, \si{}, \ang{}
\usepackage{xcolor}\usepackage{colortbl}
\usepackage{tikz}\usepackage{pgfplots}\pgfplotsset{compat=1.18}
\usepackage[siunitx,europeanresistors]{circuitikz} % circuits (élec.)
\usepackage[most]{tcolorbox}\usepackage{enumitem}\usepackage{array,booktabs}
\usepackage[a4paper,margin=2.2cm]{geometry}
\usetikzlibrary{arrows.meta,calc,angles,quotes,positioning,patterns,%
  backgrounds,shapes.geometric,decorations.pathmorphing,decorations.markings,intersections}
\definecolor{primary}{RGB}{222,49,99}\definecolor{primarydark}{RGB}{150,22,55}
\definecolor{defcol}{RGB}{0,114,178}\definecolor{methcol}{RGB}{52,92,148}
\definecolor{excol}{RGB}{0,135,90}\definecolor{attncol}{RGB}{200,40,55}
\tcbset{boxbase/.style={enhanced,breakable,boxrule=0.9pt,arc=2.5pt,
  left=3mm,right=3mm,top=2.3mm,bottom=2.3mm,fonttitle=\bfseries,coltitle=white}}
% --- boîtes TUTORIEL (noms EXACTS, ne pas renommer) ---
\newtcolorbox{cle}[1][]{boxbase,colback=defcol!6,colframe=defcol,
  title={\textsc{À retenir}\ifx\relax#1\relax\relax\else\textmd{ : \itshape #1}\fi}}
\newtcolorbox{methode}[1][]{boxbase,colback=methcol!7,colframe=methcol,sharp corners,
  title={\textsc{Méthode}\ifx\relax#1\relax\relax\else\textmd{ --- \itshape #1}\fi}}
\newtcolorbox{exemplebox}[1][]{boxbase,colback=excol!5,colframe=excol,
  title={\textsc{Exemple}\ifx\relax#1\relax\relax\else\textmd{ : \itshape #1}\fi}}
\newtcolorbox{piege}[1][]{boxbase,colback=attncol!6,colframe=attncol,
  title={\textsc{Piège}\ifx\relax#1\relax\relax\else\textmd{ : \itshape #1}\fi}}
\newtcolorbox{remarque}[1][]{boxbase,colback=black!4,colframe=black!45,
  title={\textsc{Remarque}\ifx\relax#1\relax\relax\else\textmd{ : \itshape #1}\fi}}
% --- boîtes EXERCICE / CORRIGÉ (noms EXACTS, auto-comptées) ---
\newtcolorbox[auto counter]{exo}[1][]{boxbase,colback=primary!4,colframe=primary,
  title={\textsc{Exercice}~\thetcbcounter\ifx\relax#1\relax\relax\else\textmd{ --- \itshape #1}\fi}}
\newtcolorbox[auto counter]{corrige}[1][]{boxbase,colback=excol!4,colframe=excol,
  title={\textsc{Corrigé}~\thetcbcounter\ifx\relax#1\relax\relax\else\textmd{ --- \itshape #1}\fi}}
\newcommand{\vv}[1]{\vec{#1}}\newcommand{\ds}{\displaystyle}
\newcommand{\R}{\mathbb{R}}\newcommand{\N}{\mathbb{N}}\newcommand{\Z}{\mathbb{Z}}
% (un \usepackage{fancyhdr}+en-tête est possible ; il est ignoré côté web)
% =========================================================================
\begin{document}
% THEME_TITLE: Équilibre de rotation — leviers et balances
\begin{center}
{\LARGE\bfseries\color{primarydark}Thème 5 — Équilibre de rotation : leviers et balances}\\[2pt]
{\color{primary}\itshape Statique — Fiche 3}
\end{center}
\vspace{1mm}

Un solide ne tourne pas lorsque les moments qui tendent à le faire pivoter dans un sens sont
exactement compensés par ceux qui le font pivoter dans l'autre. C'est la \textbf{deuxième} condition
d'équilibre, indépendante de $\sum\vv F=\vv 0$.

\begin{cle}[condition d'équilibre de rotation]
\[
\boxed{\;\sum_i \mathcal{M}_\Delta(\vv F_i) = 0\;}
\qquad\Longleftrightarrow\qquad
\sum \mathcal{M}_{\text{horaire}} = \sum \mathcal{M}_{\text{anti-horaire}}
\]
La somme est \textbf{algébrique} (chaque moment avec son signe). Le moment « moteur » est équilibré
par la somme des moments « résistants ».
\end{cle}

\begin{cle}[le levier]
Pour un levier à pivot, deux poids s'équilibrent quand leurs moments sont égaux :
\[
\boxed{\;G_A\times d_A = G_B\times d_B\;}
\]
($d_A$, $d_B$ = distances au pivot). Conséquence : un \emph{petit} poids placé \emph{loin} du pivot
équilibre un \emph{grand} poids placé \emph{près} — c'est tout le principe du levier.
\end{cle}

\begin{methode}[résoudre un problème de statique complet]
\begin{enumerate}[leftmargin=1.6em,topsep=1pt,itemsep=1pt]
  \item Faire le \textbf{bilan des forces} (poids, réactions d'appui, tensions\dots).
  \item Choisir un \textbf{axe} : le placer sur un appui \emph{annule le moment de sa réaction}
        (bras nul) et élimine une inconnue.
  \item Écrire l'\textbf{équilibre de rotation} $\sum\mathcal{M}=0$ autour de cet axe.
  \item Écrire l'\textbf{équilibre de translation} $\sum\vv F=\vv 0$.
  \item Résoudre le système.
\end{enumerate}
\end{methode}

\begin{cle}[poutre / table sur deux appuis]
Deux appuis $A$ et $B$ exercent des réactions $R_A$ et $R_B$. On dispose de deux équations :
$\;R_A+R_B = \text{poids total}\;$ (translation) et $\;\sum\mathcal{M}=0\;$ (rotation autour d'un
appui). Le résultat clé : \textbf{l'appui le plus proche de la charge en supporte la plus grande part}.
\end{cle}

\begin{exemplebox}[levier à deux bras]
$G_A=\qty{20}{N}$ à $d_A=\qty{0.3}{m}$ du pivot ; à quelle distance $d_B$ placer $G_B=\qty{30}{N}$
pour l'équilibre ? $G_A d_A = G_B d_B \Rightarrow d_B = \dfrac{20\times 0{,}3}{30}=\qty{0.2}{m}$. Le
poids le plus grand se place plus \emph{près} du pivot.
\end{exemplebox}

\begin{piege}[il faut les deux conditions]
$\sum\vv F=\vv 0$ \emph{seul} ne suffit pas (un couple le vérifie et fait pourtant tourner) ;
$\sum\mathcal{M}=0$ \emph{seul} non plus. L'\textbf{équilibre complet} exige les \emph{deux} en même temps.
\end{piege}

\end{document}
